\section*{ Introduction }
\addcontentsline{toc}{section}{Introduction}
\markboth{INTRODUCTION}{INTRODUCTION}

Le transport de bagages par des véhicules guidés automatisés (AGV) en environnement aéroportuaire soulève des problématiques complexes de navigation autonome, de coordination entre agents et d'évitement d'obstacles dans des environnements dynamiques. Afin d'étudier ces problématiques, un simulateur a été développé pour modéliser différents scénarios de navigation et expérimenter des approches d'apprentissage par renforcement dans un cadre contrôlé et reproductible.

\vspace{10pt}

Dans une première étape, le simulateur est utilisé dans un contexte plus général de navigation de robots mobiles autonomes (AMR). Il permet de créer des environnements configurables composés d'agents, d'obstacles statiques et d'obstacles dynamiques, puis d'entraîner et d'évaluer des politiques de navigation. La version actuelle du simulateur repose sur l'algorithme Soft Actor-Critic (SAC) et se concentre sur l'apprentissage d'un agent unique dans des environnements de difficulté croissante.

\vspace{10pt}

Ce manuel a pour objectif de présenter l'utilisation du simulateur et d'accompagner l'utilisateur dans sa prise en main. Il décrit successivement les étapes d'installation, l'organisation générale du projet, les fichiers de configuration, les principaux modes d'exécution ainsi que les outils permettant de visualiser et d'analyser les résultats générés. Il s'adresse aussi bien aux utilisateurs souhaitant reproduire les expériences existantes qu'aux développeurs désirant créer de nouveaux scénarios ou adapter le simulateur à leurs propres besoins.
